\section{A bare primitive on the justified ramification set}\label{sec:bareprimitive}
Fix one of the nine prime beta supports $S$ in Section~\ref{sec:patterns}. Let $\Sigma$ consist of $S$, the places above $2,3,5,7$, and the two real places. The residual classes factor through $\Gsig$. At every beta support place, $\chi$ is trivial modulo seven and the natural local class is zero. Hence its normalized cocycle is identically zero there. At the fixed places, beta is a local coboundary under the inherited contract.

There are three questions to keep distinct:
\begin{enumerate}
\item Does the bare restricted equation $dp_0=-\alpha\cup\beta$ have a solution?
\item Does one such global solution satisfy all prescribed first trace and toric conditions?
\item Does a matching first solution extend to the upper and higher integral geometric stages?
\end{enumerate}
The first question can now be answered using finite set duality. Its answer does not assume an actual higher representation in order to construct the first primitive.

\begin{lemma}[The dual Kummer kernel]\label{lem:sha1}
For every finite $\Sigma$ containing the prime above seven and all real places,
\[
\Sha^1_\Sigma(K,\mu_7)
=\ker\left(H^1(\Gsig,\mu_7)
\longrightarrow\prod_{v\in\Sigma}H^1(K_v,\mu_7)\right)=0.
\]
\end{lemma}
\begin{proof}
The discriminant of $K$ is five. Minkowski's ideal bound is $\sqrt5/2<2$, so every ideal class has a representative of norm one. Thus $\Cl(K)=1$ \cite{milneANT2020}.

Its units are $\{\pm\phi^e:e\in\Z\}$. For an elementary verification, scale a unit by a power of $\phi$ and by a sign so that its first real embedding lies in $[1,\phi)$. If its value $a$ were strictly between the endpoints, a norm minus one unit would have integral trace $a-a^{-1}$ strictly between zero and one. A norm plus one unit would have trace $a+a^{-1}$ strictly between two and $\sqrt5<3$. Both are impossible. The normalized unit is therefore one, proving the assertion.

A Kummer class in the restricted group has valuations divisible by seven outside $\Sigma$. Class number one lets us remove those valuations by a global seventh power and leave a $\Sigma$ unit. If its local class is zero at every place in $\Sigma$, the remaining prime valuations are also divisible by seven. Remove these as well. Since minus one is its own seventh power, the class is represented by $\phi^e$.

At $F=\Q_7(\sqrt5)$, exact arithmetic modulo $49$ gives
\begin{equation}\label{eq:phi-nonseventh}
\frac{\phi^{48}-1}{7}=5\overline{\sqrt5}\ne0\quad\text{in }\F_{49}.
\end{equation}
If a unit is a seventh power, its forty eighth power is one modulo forty nine. Indeed, any unit has forty eighth power in $1+7\Oloc_F$, and the seventh power of this group is contained in $1+49\Oloc_F$. Formula \eqref{eq:phi-nonseventh} shows that $\phi$ is not a seventh power in $F$.

If $e$ is nonzero modulo seven and $\phi^e$ is a seventh power, a Bezout identity between $e$ and seven makes $\phi$ a seventh power too. Therefore local triviality forces $e\equiv0\pmod7$, and the global Kummer class is zero.
\end{proof}
The lemma is uniform in finite $\Sigma$. It does not use a prime sample to test arbitrary global units. The single seven adic symbol is sufficient only after the class number and unit reductions have removed every other possible coordinate.

\begin{theorem}[Conditional restricted bare primitive]\label{thm:bareprimitive}
Assume the stated residual and local properties for $\alpha\in Z^1(\Gsig,k(\chi))$ and $\beta\in Z^1(\Gsig,k(\chi^{-1}))$. Then there is a normalized continuous cochain $p_0:\Gsig\to k$ satisfying
\[
dp_0=-\alpha\cup\beta.
\]
No auxiliary ramification outside $\Sigma$ is needed.
\end{theorem}
\begin{proof}
The Cartier dual of trivial $\F_7$ is $\mu_7$. The nondegenerate $\F_7$ trace pairing on $k$ identifies the dual of trivial $k$ with $k(\chi)$. Lemma~\ref{lem:sha1} in two $\F_7$ coordinates therefore gives $\Sha^1_\Sigma(K,k(\chi))=0$.

Finite set Poitou--Tate duality pairs this group perfectly with $\Sha^2_\Sigma(K,k)$. Every prime dividing the module order belongs to $\Sigma$, as required by the cited theorem. Odd primary real Tate terms vanish. Thus localization
\begin{equation}\label{eq:restrictedH2}
H^2(\Gsig,k)\longrightarrow\bigoplus_{v\in\Sigma}H^2(K_v,k)
\end{equation}
is injective \cite{milneADT2006}.

At beta support, $\alpha$ is identically zero locally, so its cup with beta is zero. At a fixed place, beta is a coboundary, so the local cup class is zero. The real terms vanish. Injectivity of \eqref{eq:restrictedH2} gives $[\alpha\cup\beta]=0$ globally in the restricted group. A continuous one cochain primitive exists by the definition of cohomological vanishing. Evaluating the primitive equation at the identity gives $p_0(1)=0$.
\end{proof}
It would be invalid to replace \eqref{eq:restrictedH2} by an assertion about unrestricted absolute cohomology and infer the same ramification control. The dual kernel has been killed inside the restricted problem itself.

For the frozen supports, each supporting local field has residue norm one modulo seven and contains $\mu_7$. Local duality gives one $k$ dimensional $H^2$ term there. The four fixed fields do not contain $\mu_7$ and have zero $H^2(K_v,k)$. Global dual invariants also vanish. The final portion of the Poitou--Tate sequence then identifies the restricted $H^2$ with $k^9$. The dimension nine specializes the frozen support, whereas Lemma~\ref{lem:sha1} held for arbitrary finite $\Sigma$.

The theorem proves existence of an origin in the bare primitive torsor. It does not compute values of $p_0$ on labelled Frobenius or inertia, a defining field for $p_0$, or its affine compatibility class. Those numerical values are exactly what the final evaluator still lacks.

\subsection{The unique local fixed trace primitive}
At a fixed place $v$ above $2,3,5,7$, write
\[
\beta(g)=(\chi(g)^{-1}-1)\lambda_v.
\]
The local separating element makes $\lambda_v$ unique. Define
\begin{equation}\label{eq:pfix}
p_v^{\mathrm{fix}}(g)=\lambda_v\chi(g)^{-1}\alpha(g).
\end{equation}
Direct cocycle expansion yields
\begin{align}
dp_v^{\mathrm{fix}}(g,h)&=-\alpha(g)\chi(g)^{-1}\beta(h),\label{eq:pfixdiff}\\
(\chi(g)-1)p_v^{\mathrm{fix}}(g)+\alpha(g)\beta(g)&=0.\label{eq:pfixtrace}
\end{align}
The second equation is precisely the first trace condition in \eqref{eq:adjointtrace}.

If another local primitive satisfies the same trace condition, their difference $z$ is an ordinary character and $(\chi(g)-1)z(g)=0$ for every $g$. If $\chi(g)\ne1$, this already gives $z(g)=0$. For $\chi(g)=1$, choose $h$ with $\chi(h)\ne1$. Both $z(h)$ and $z(hg)$ vanish, so additivity gives $z(g)=0$. Thus \eqref{eq:pfix} is the unique local cochain satisfying both equations in this fixed frame. It is not uniqueness of a global compatible primitive.

For a bare global $p_0$, the difference $p_v^{\mathrm{fix}}-p_0|_v$ is an ordinary local character. Only that difference has the class field coordinates used below. Neither inhomogeneous cochain can be treated individually as an ordinary character.
